\setcounter{chapter}{5}
\chapter{Phase 1C: Production Deployment}
\minitoc
\graphicspath{{Chapter6/figures/}}

%==============================================================================
\pagestyle{fancy}
\fancyhf{}
\fancyhead[R]{\bfseries\rightmark}
\fancyfoot[R]{\thepage}
\renewcommand{\headrulewidth}{0.5pt}
\renewcommand{\footrulewidth}{0pt}
\renewcommand{\chaptermark}[1]{\markboth{{\chaptername~\thechapter.\ #1}}{}}
\renewcommand{\sectionmark}[1]{\markright{\thechapter.\thesection~ #1}}

\begin{spacing}{1.2}
%==============================================================================

\section*{Introduction}
% TODO:
%   - Frame Phase 1C as the operational realisation of the system
%     topology designed in Chapter II Section 5.
%   - Engineer's posture: what to monitor, what to back up, what
%     SLO to target, how to release safely. Tools are illustrations,
%     not subject matter.
%   - State the milestone closed at the end of Phase 1C: live in
%     production at \texttt{clmpilot.com}, monitored, backed up,
%     with a repeatable release path.

%--------------------------------------------------------------------
\section{Production topology}
% TODO: One diagram of what actually runs in production. Differences
% from the topology diagram in Chapter II (if any) explained briefly.
%   - Single Hetzner CX22 host, Docker Compose stack.
%   - Services: backend, frontend, intelligence stub, postgres, redis.
%   - Reverse proxy: nginx on the host (not inside the compose stack).
%   - Outbound dependencies: SendGrid, Google Drive, SharePoint.
% AVOID: ``how I provisioned the server'' tutorial. The provisioning
% recipe lives in Appendix A.3.

\begin{figure}[!ht]\centering
  % \includegraphics[scale=0.7]{production-topology.png}
  \caption{Production deployment topology.}
  \label{fig:prod-topology}
\end{figure}

%--------------------------------------------------------------------
\section{Infrastructure choices and trade-offs}
% TODO: Engineer's decisions, justified.
%   - Hetzner CX22 vs. AWS/GCP/Azure: cost profile for early-stage
%     CLM, EU data residency, simple billing.
%   - Docker Compose vs. Kubernetes for the MVP: K8s is over-engineered
%     for a single-host MVP; the path to K8s is on the Phase 3 roadmap.
%   - Apex-only domain (\texttt{clmpilot.com}, no \texttt{app.} subdomain)
%     -- a small architectural decision that matters operationally
%     ([[project-apex-domain]]).
%   - Email via SendGrid HTTP API rather than self-hosted SMTP -- removes
%     deliverability and DNS-warming burden ([[project-smtp-provider]]).
\dots

%--------------------------------------------------------------------
\section{Edge and transport security}
% TODO: At the security-posture level, not the nginx config level.
%   - TLS 1.3 with Let's Encrypt, automated renewal.
%   - HTTP-to-HTTPS redirection enforced at the proxy.
%   - SSH hardening: public-key only, no root password auth.
%   - One short paragraph naming the controls; the nginx and
%     systemd files live in Appendix A.3.
\dots

%--------------------------------------------------------------------
\section{Observability and operational posture}
% TODO: What the team can see and react to in production.
%   - Health checks at the compose level (every container has one).
%   - Application-level error reporting to email (MVP-lite: SendGrid
%     alerts instead of a paid APM).
%   - Backup strategy: daily Postgres dumps with off-host retention,
%     restore procedure tested on schedule.
%   - SLO target for the MVP: not a 99.99\% promise, but a stated
%     posture (~99\% within business hours; planned-maintenance
%     windows).
%   - Engineer's note: the auditability built in at the event-store
%     level is itself an observability lever -- any production
%     anomaly can be replayed from events.
\dots

%--------------------------------------------------------------------
\section{Release process and MVP validation}
% TODO: How a change goes from develop branch to production.
%   - The release pipeline: develop -> main PR -> image build -> tag bump
%     in the deployment repo -> \texttt{deploy.sh} on the host.
%   - The lockstep behaviour of \texttt{deploy.sh}: writes the same
%     version to all three component repos. Documented as an
%     architectural property to be aware of; the per-component bypass
%     procedure is in the runbook ([[feedback-deploy-script-lockstep-pain]]).
%   - End-to-end validation gate (Playwright suite) before promotion.
%   - The MVP launch checklist as a release-process artefact, not a
%     to-do list -- prove each checklist item exists as an automated
%     check.
\dots

%--------------------------------------------------------------------
\section*{Conclusion}
% TODO: Recap (production deployment realised on the topology of
% Chapter II, validated by a Phase 1A/1B/1C sequence that culminated
% in the live MVP at \texttt{clmpilot.com}). Bridge: the General
% Conclusion (and the deferred Phase 2/3 chapters) extend the story.
\dots

%==============================================================================
\end{spacing}
